% Build in this directory: pdflatex -interaction=nonstopmode -halt-on-error answer_exercise1.tex
% Exercise statements from class1PT26.tex.
\documentclass[a4paper,10pt,pdflatex,ja=standard,jafont=ipaex-type1,
  magstyle=real,scale=1]{bxjsarticle}
\usepackage{lmodern}
\usepackage{amsmath,amssymb,mathtools}
\usepackage{geometry}
\geometry{reset,left=22mm,right=22mm,top=20mm,bottom=17mm,
  headheight=13pt,headsep=6mm,footskip=10mm}
\usepackage{xcolor}
\usepackage{fancyhdr}
\usepackage{hyperref}

\definecolor{ink}{HTML}{193A5C}
\definecolor{accent}{HTML}{237D82}
\definecolor{soft}{HTML}{EAF3F3}
\newcommand{\coursename}{確率論と統計学 I}
\newcommand{\F}{\mathcal F}
\hypersetup{unicode=true,colorlinks=true,urlcolor=accent,linkcolor=ink,
  pdfauthor={中島秀太},pdftitle={第1回：演習問題と解答}}
\pagestyle{fancy}
\fancyhf{}
\fancyhead[L]{\small\color{ink}\coursename{}・演習問題と解答}
\fancyhead[R]{\small\color{ink}第1回}
\fancyfoot[C]{\small\color{ink}\thepage}
\renewcommand{\headrulewidth}{0.3pt}
\setlength{\parindent}{0pt}
\setlength{\parskip}{2pt}
\setlength{\abovedisplayskip}{4pt}
\setlength{\belowdisplayskip}{4pt}
\setlength{\abovedisplayshortskip}{3pt}
\setlength{\belowdisplayshortskip}{4pt}
\newcommand{\topic}[1]{\par\vspace{6pt}
  {\large\sffamily\mdseries\color{accent}#1}\par\vspace{3pt}}
\newcommand{\handouttitle}[1]{%
  \begin{center}
    {\Large\color{ink}#1}\par\vspace{3mm}
    {\color{accent}\rule{0.88\linewidth}{0.7pt}}
  \end{center}\vspace{-2mm}}
\newcommand{\problem}[1]{\par\vspace{3pt}\begingroup
  \setlength{\fboxsep}{5pt}
  \colorbox{soft}{\parbox{\dimexpr\linewidth-2\fboxsep\relax}{#1}}%
  \endgroup\par\vspace{4pt}}
\newcommand{\answerlabel}{{\sffamily\mdseries\color{ink}解答}\par}

\begin{document}
\fontsize{10pt}{13.5pt}\selectfont
\handouttitle{第1回\quad 演習問題と解答}

\topic{演習1}
\problem{公平なサイコロを独立に2回投げる試行の離散確率空間
$(\Omega,\F,P)$ を構成し、出た目の合計が3となる確率を求めよ。}

\answerlabel
1回目の出目を $i$、2回目の出目を $j$ として、
$$
  \Omega=\{1,2,3,4,5,6\}^2
    =\{(i,j):1\le i,j\le6\},\qquad \F=2^\Omega
$$
とおく。公平性と独立性より、各結果の重みは
$$
  p_{(i,j)}=\frac16\cdot\frac16=\frac1{36},\qquad
  \sum_{(i,j)\in\Omega}p_{(i,j)}=36\cdot\frac1{36}=1.
$$
したがって、$|A|$ を集合 $A$ の要素数として、
$$
  P(A)=\sum_{(i,j)\in A}p_{(i,j)}=\frac{|A|}{36}
  \qquad(A\in\F)
$$
と定める。$P(\Omega)=1$ であり、互いに素な事象列 $(A_n)_{n\ge1}$ について
$$
  P\!\left(\bigcup_{n=1}^\infty A_n\right)
  =\frac1{36}\sum_{n=1}^\infty|A_n|
  =\sum_{n=1}^\infty P(A_n)
$$
が成り立つので、$(\Omega,\F,P)$ は求める離散確率空間である。
合計が3となる事象は $S=\{(i,j)\in\Omega:i+j=3\}=\{(1,2),(2,1)\}$ だから、
$$
  \boxed{P(S)=\frac2{36}=\frac1{18}}.
$$

\topic{演習2}
\problem{$(\Omega,\F,P)$ を確率空間とする。任意の $A,B\in\F$ に対して、
$A\cap B\in\F$ および $P(A\cup B)=P(A)+P(B)-P(A\cap B)$ を示せよ。}

\answerlabel
$A^c=\Omega\setminus A$ とする。$\sigma$-加法族は補集合と有限和について閉じているので、
ド・モルガンの法則より
$$
  A\cap B=(A^c\cup B^c)^c\in\F.
$$
また、$\varnothing=\Omega^c\in\F$ である。可算加法性に
$A_1=\Omega$, $A_n=\varnothing$ $(n\ge2)$ を代入すると、
$1=1+\sum_{n=2}^\infty P(\varnothing)$ より $P(\varnothing)=0$。
したがって、可算加法性で残りの事象を空集合とすれば有限加法性が得られる。

$$
  C=A\setminus B=A\cap B^c,\qquad
  D=B\setminus A=B\cap A^c,\qquad E=A\cap B
$$
とおく。いずれも $\F$ に属し、互いに素である。
$\sqcup$ は互いに素な集合の和を表すものとすると、
$$
  A=C\sqcup E,\qquad B=D\sqcup E,\qquad
  A\cup B=C\sqcup D\sqcup E.
$$
有限加法性から
$$
  \begin{aligned}
    P(A)+P(B)-P(A\cap B)
      &=\bigl(P(C)+P(E)\bigr)+\bigl(P(D)+P(E)\bigr)-P(E)\\
      &=P(C)+P(D)+P(E)\\
      &=P(A\cup B).
  \end{aligned}
$$
よって、$\boxed{P(A\cup B)=P(A)+P(B)-P(A\cap B)}$ が成り立つ。

\end{document}
